\subsection{3. Mục tiêu đề tài}
\indent\indent Mục tiêu tổng quát của đề tài là nghiên cứu và xây dựng hệ thống phát hiện tài khoản giả mạo (Sybil) trên mạng xã hội phi tập trung (Web3) bằng việc tận dụng mô hình đồ thị trí tuệ nhận thức ràng buộc kết hợp với kiến trúc phân lớp tách rời (Decoupled ML). Đề tài hướng tới việc giải quyết bài toán thiếu hụt nhãn dữ liệu thực tế và nâng cao hiệu năng phân lớp thông qua việc học các mối quan hệ phức tạp từ hành vi tương tác, nội dung ngữ nghĩa và các ràng buộc tài sản trên chuỗi khối. Kết quả nghiên cứu sẽ cung cấp cơ sở thực nghiệm cho việc ứng dụng mạng nơ-ron đồ thị (GNN) vào bài toán bảo mật hệ sinh thái SocialFi. Các mục tiêu cụ thể mà đề tài hướng đến bao gồm:

\begin{itemize}[label={--}, itemsep=1pt]
    \item \textbf{Thu thập và tiền xử lý dữ liệu:} Trích xuất dữ liệu thực tế từ mạng xã hội Lens Protocol thông qua nền tảng Google BigQuery với ba phạm vi thời gian (tuần, tháng, quý). Xây dựng không gian đặc trưng 396 chiều cho mỗi nút, bao gồm: ngữ nghĩa được trích xuất từ cấu trúc hồ sơ bằng mô hình Sentence-BERT và chỉ số thống kê định lượng về hành vi, thiết lập bộ dữ liệu đầu vào tiêu chuẩn cho mô hình học sâu.
    
    \item \textbf{Xây dựng đồ thị đa tầng:} Mô hình hóa dữ liệu thành một đồ thị đa quan hệ với 4 tầng kiến trúc chuyên biệt: Theo dõi, Tương tác, Tương đồng hồ sơ và Đồng sở hữu ví. Tích hợp cơ chế tinh chỉnh trọng số cơ sở cho từng tầng nhằm định hướng mô hình tập trung vào các ràng buộc tài chính on-chain và dấu vết hành vi bầy đàn cốt lõi.
    
    \item \textbf{Học không giám sát và tự động gán nhãn:} Khắc phục bài toán thiếu nhãn thực tế thông qua một quy trình tự động hóa. Ứng dụng bộ mã hóa tự động đồ thị (GAE) kết hợp kiến trúc GATv2 để học biểu diễn nút, tiếp nối bằng thuật toán phân cụm K-Means và hệ thống luật chuyên gia để gán nhãn giả lập phân loại 4 mức độ rủi ro.
    
    \item \textbf{Tối ưu hóa phân lớp bằng kiến trúc tách rời (Decoupled ML):} Thực hiện nghiên cứu loại trừ toàn diện trên 20 tổ hợp mô hình, kết hợp giữa 5 kiến trúc GNN tiên tiến (GAT, GCN, PNA,...) và 4 thuật toán máy học truyền thống (Random Forest, XGBoost,...). Mục tiêu nhằm vượt qua giới hạn "hộp đen" và đạt được điểm số Macro F1 vượt trội.
    
    \item \textbf{Phát triển và triển khai hệ thống ứng dụng thực tiễn:} Đóng gói mô hình thành một ứng dụng web hoàn chỉnh (sử dụng Next.js, FastAPI, Modal). Hệ thống cung cấp tính năng tra cứu điểm rủi ro qua định danh tài khoản, hoặc thực thi toàn bộ luồng trên một khoảng thời gian tùy chọn. Giao diện tích hợp khả năng trực quan hóa đồ thị cục bộ và biểu đồ diễn giải nguyên nhân rủi ro, đảm bảo tính minh bạch cho người dùng cuối.
\end{itemize}